% !TeX root = ../papers/nc-p3.tex
\section{楕円四次点成分の二次核と全次数の判定}
\label{sec:elliptic-point-quartic}
前節の有限双対Frobenius窓のうち、点グラフの一次元部分に滑らかな楕円四次曲線 $E$ を持つものを取る。
平面内の滑らかな四次曲線は種数3なので $E$ は非退化である。従って
$V_i=H^0(E,L_i)$ は完全線型系である。$D_{ij}=H^0(E,L_i\cdots L_{j-1})$ ($i\le j$)
を積で結び、$A$ を元の窓の二次代数とする。定理\ref{thm:elliptic-pencil-all-as}の
評価核のフィルトレーションは各 $L_i$ の次数4だけを使うため、ここでも $D$ は一次生成かつ二次関係である。
$A\to D$ の核を $I$、$W_i=I_{i,i+2}$ と書く。$\dim W_i=2$ である。

\begin{theorem}\label{thm:elliptic-quartic-as-test}
$A$ が標準Koszul左右AS正則となる必要十分条件は、全四位置で
\[
 \ker(W_i\otimes W_{i+2}\longrightarrow A_{i,i+4})
\]
が階数2のテンソルを含むことである。そのとき核はその一本の直線であり、
$W$ の生成する部分代数は偶奇それぞれで次数を倍にした二変数多項式型の
添字付き代数である。$A$ はその上で左右自由である。
\end{theorem}
\begin{proof}[切断環を係数とするホモロジー]
始点 $i$ の右Koszul候補複体に $D$ をテンソルし終点 $i+n$ を取る。
これは前節の分裂局所完全な束複体に共通の直線束を掛けて切断を取ったもので、
項 $j$ の偏極は $L_{i+j}\cdots L_{i+n-1}$（逆順の区間は逆束）である。
$n\ge5$ なら全項の $H^1$ が零なので完全である。$n=4$ では $H^1$ は項4の
$H^1(\mathcal O_E)=k$ だけである。二行のhypercohomology系列は全体が零へ収束し、
この一項からの二階微分が切断列の項2のホモロジーへの同型となる。
$n=3$ で $H^1$ を持つ二項の間の射は同型である。その双対は末尾文字の完全線型系
$V_{i+3}\to H^0(L_{i+3})$ そのものであり、三次の係数との同定はFrobenius積対合による。
$n=2$ では六関係の $V_i\otimes V_{i+1}$ への包含と全切断乗法を使い、
項1のみに二次元が残る。$n=1,0$ は生成元の同型と次数零の単純を与える。
従って得られた正ホモロジーは、項1の次数2に二個、項2の次数4に一個である。
反対向きにも同じ議論が通る。

元の単純加群の極小自由解消は最初の二微分までこの候補複体と一致する。
三次の四つの重なりは独立な最小syzygyであり、極小解消の項3にこれらを含められる。
従って実際の $\operatorname{Tor}_1^A(k,D)$ は上の項1に等しく、
$\operatorname{Tor}_2^A(k,D)$ は上の項2の商である。左右を交換すれば、
各 $e_iD$ の右表示で $I_{i,*}=W_i A$、その関係加群は次数4の高々一元で生成される。
左表示でも $I_{*,j}=A W_{j-2}$ である。$A$ が標準なら、候補複体は最初から解消なので、
$D$ の極小左右解消はともに長さ2、生成次数は $0,2,4$、個数は $1,2,1$ である。
\end{proof}

\begin{proof}[十分性]
仮定のテンソル $r_i$ は $W_i\otimes A_{i+2,i+4}$ においても非零であり、
上の唯一のsyzygyを生成する。矢 $W_i$ と関係 $kr_i$ で定まる抽象代数を $P$ とする。
偶奇ごとに一位置の基底を決め、両方向に次の基底を選べば、
非退化な関係をいつも $x y-y x$ にできる。従って $P$ の長さ $2h$ 成分の次元は $h+1$ である。
$D$ の各成分の基底を $A$ に持ち上げ、左から $P$ の単項式を掛けると $A$ を張る。
これは $I_{i,*}=W_i A$ から次数帰納で従う。右からも張る。
従って各成分は $H_D(z)/(1-z^2)^2=(1-z)^{-4}$ の係数以下の次元を持つ。

逆向きの不等式には終点 $j$ を固定し、$a_n=\dim A_{j-n,j}$ とする。
上記の右表示から $a_n-2a_{n-2}+a_{n-4}=\dim D_{j-n,j}+k_n$、
ここで $k_n\ge0$ は最後の自由加群からの写像の核次元である。
$(1-z^2)^{-2}$ の係数が正または零なので両不等式は等号、全 $k_n=0$ である。
張るために使った基底は左右とも独立で、$P$ は $A$ に埋まり、$A$ は左右 $P$ 自由である。
$e_i D$ は $P$ の始点単純を $A$ に誘導した加群である。自由性によりその解消も誘導でき、
随伴と $P$ の二変数Koszul解消から $\underline{\operatorname{Ext}}_A^q(e_iD,A)=0$ ($q<2$) を得る。
ここで係数の $A$ は右 $P$ 自由でもある。環変更により $A$ の右単純から $A$ への
次数0,1のExtも消える。左側でも同様である。
標準候補複体の先頭二写像は生成・二次表示で完全であり、末尾二箇所も
反対側の先頭二写像の双対であるためこのExt消滅で完全となる。
残る中央は既に求めたHilbert関数とEuler式で完全である。
Frobenius対合による左右解消の双対性がAS条件を与える。
\end{proof}

\begin{proof}[必要性]
標準の場合を考える。唯一の楕円四次成分は標準四移動で保たれるので、
$A\to D$ は四頂点の正次数付き代数の全射に折り返せる。以下の双対 $(-)'$ は局所有限な次数付き双対である。
$D$ の長さ2の極小右解消を双対化して得る
$\underline{\operatorname{Ext}}_A^2(D,A)$ が左から $I$ で零化されることを示す。

まず有限次元右加群 $F$ には自然な左右双対
$\underline{\operatorname{Ext}}_A^4(F,A)=F'$（頂点・次数の標準移動とNakayamaねじれを除く）があり、
他のExtは零である。この同定の作用も確認しておく。$T_j$ を次数 $j$ のKoszul重なりとし、
対角解消 $A\otimes_{A_0}T_j\otimes_{A_0}A$ の微分を最初と最後の文字の積の符号付き和とする。
これは片端の次数でフィルターして単純のKoszul解消を使えば対角を解消する。
左から $F$ をテンソルした右自由解消の $A$ 双対の項は
$A\otimes_{A_0}T_j^*\otimes_{A_0} F'$ である。
$T_j^*$ をFrobenius積対合で相補位置の $T_{4-j}$ と同定すると、両端の微分は
逆順の二つの文字の積に移る。一周する端でだけtraceの巡回によるNakayama移動が入る。
従ってこれは $F'$ の左対角解消を逆に読んだものである。符号の次数自己同型と標準四移動は $I$ を保つ。

$F_N=A/A_{\ge N}$ とし、右加群の次数付きtorsion関手を
$\Gamma_A(M)=\varinjlim_N\underline{\operatorname{Hom}}_A(F_N,M)$ と書く。
$F_N$ の左作用を前合成すると結果は右加群になる。上の自然な同定と
$\varinjlim F_N'=A'$ により、$R\Gamma_A(A)$ は次数4だけに、標準のねじれ・移動を持つ
双加群 $A'$ を持つ。右 $D$ の有限自由右 $A$ 解消にこの関手を施せば、各自由項の導来torsionは
その $A$ 双対の次数付き双対にこの同じ移動を施したものである。
自由項間の射は左からの行列なので上の双加群同定の自然性で微分も対応する。
特に $\underline{\operatorname{Ext}}_A^2(D,A)$ は $R^2\Gamma_A(D)$ の次数付き双対のねじれ・移動である。

$R\Gamma_A(D)$ の右作用が $D$ を通ることは、ここではnoetherian仮定なしに確認できる。
右 $D$ 入射加群 $J$ に対して環変更で
$\underline{\operatorname{Ext}}_A^q(F_N,J)
=\underline{\operatorname{Hom}}_D(\operatorname{Tor}_q^A(F_N,D),J)$ となる
（次数移動した入射も使う）。$D$ の左 $A$ 解消の生成次数は $0,2,4$ であるから、
$q>0$ のTorは内部次数 $[N,N+3]$ にしかない。下限は $N$ 未満で元の解消のまま、
上限は $F_N$ と最も高い生成次数から従う。よって $N'\ge N+4$ からの遷移は零であり
$J$ は $\Gamma_A$ 非輪状である。$D$ 入射解消で計算すれば
$\Gamma_A(J)=\Gamma_D(J)$ は右 $D$ 加群であり、比較は右 $A$ 作用を保つ。
従って要求した $I$ による零化が従う。

最後に各始点の $D$ の右解消の唯一の最終列を $(u_i,v_i)^t$ とする。
$u_i,v_i\in A_{i+2,i+4}$ であり、その双対の余核は左巡回加群である。
その最小生成元に $W_{i+2}\subset I$ を左から作用させると零になる。
この内部次数で双対微分の像は $u_i,v_i$ のスカラー線型包に等しい。
よって $W_{i+2}=\langle u_i,v_i\rangle$、二係数は独立で、syzygyは $W_i\otimes W_{i+2}$ 内の階数2である。
\end{proof}

\begin{corollary}\label{cor:elliptic-point-finite-atlas}
この枝の添字保持同型類は次の有限判定付きパラメータと有限群同一視で尽くされる。
楕円曲線の各 $j$ について原点付きWeierstrass代表 $(E,o)$ を固定し、
$L_0=\mathcal O(4o)$、$L_1\in\operatorname{Pic}^4(E)$、$K\in E[12]$、
$L_2=K L_1^2/L_0$ とする。$a=0,1$ で切断乗法の八次元核から
$R_a\in\operatorname{Gr}(6,8)$ を一つずつ取る。
\begin{enumerate}
\item 二窓の三次重なり $U$ の次元が4であることを課し、$V_3=U^*$ と包含テンソルから
定理\ref{thm:frobenius-tensor-reconstruction}の条件と $R_a(w)=R_a$ ($a=0,1$) を検査する。
\item 四つの点対応の左右評価で階数2以下の点がないことと、最初の像が一次元であることを検査する。
各 $W_i$ は $E$ 上に制限して零になる二文字式を $R_i$ で割って得る。
復元した二辺への $E$ の埋込みは核直線の射で運ぶ。
\item 定理\ref{thm:elliptic-quartic-as-test}の四つの乗法核を検査する。
\end{enumerate}
生き残るパラメータは同じ $j$ で $E[4]\rtimes\operatorname{Aut}(E,o)$ の引戻しにより
二つの束と二つの関係が全て対応する時かつその時に限り同型となる。
最小周期は同じ比較を一移動、二移動した二窓に施して決まる。
\end{corollary}
\begin{proof}
任意の次数4の偏極は原点の平行移動で指定の $L_0$ にでき、残る平行移動は $E[4]$ である。
切断乗法の核の条件により最初の二辺は指定の $E$ のグラフを含む。
二番目の検査の下では全辺の射影像が同型で、他の成分や埋込み点を許しても
定理\ref{thm:point-graph-curve-degree}で楕円四次成分は一意となる。
三番目の検査は必要十分で、全ては有限次元の線型代数と射影消去で行える。
同型は唯一の曲線を各位置で保ち、グラフにより各位置の曲線同型は共通である。
完全線型系なのでその引戻しが文字の写像をスカラー倍まで決める。
逆に二窓を対応させる写像は系\ref{cor:two-window-extension}により全体に延びる。
移動した二窓は最初の偏極を再び $\mathcal O(4o)$ に直して同じ比較を行う。
\end{proof}
